\documentclass[10pt]{article}
\usepackage{icmlko}

\icmlmeta{IP 파운데이션 월드모델}{23}{값진 축을 실제로 확보했다}{노트 22의 예측이 맞았다 --- 매장 노출도를 33\%에서 90\%로 올리니 아이돌 전이가 전부 강해졌다}

\begin{document}

\twocolumn[
\icmlheader

\begin{abstract}
노트 22는 매장 노출도를 최우선 확보 대상으로 지목했다 --- 팝업에서 빼면 전이가
43\% 악화되는 가장 값진 축인데 아이돌에서 32\%, 게임에서 0\%만 관측된다. 확보를
시도했다. \textbf{두 번 실패하고 세 번째에 됐다.} 노트 20의 교훈대로 이름이 아니라
물리량으로 정의하려고 ``경쟁 밀도의 역수''를 잡았는데, Steam은 검색 API의 날짜
필터가 작동하지 않아(모든 구간에서 116,832 반환) 월별 출시 수를 못 얻었고, 아이돌은
표본 편향 때문에 무효였다(월별 데뷔 수와 초동의 상관 $+$0.078, 부호도 예상과 반대).
세 번째로 \textbf{소속사 사전 데뷔 건수}를 썼다 --- 라벨이 필요 없고 커버리지가
33\%에서 \textbf{90\%}로 오르며 단조 관계가 뚜렷하다(첫 데뷔 소속사 log10 4.70,
1건 경험 5.22, 2건 이상 5.34). 결과는 노트 22의 예측대로다. \textbf{아이돌이 낀
세 방향이 모두 16--26\% 강해졌다} --- 팝업 $\to$ 아이돌 $-$0.0716에서 $-$0.0905,
아이돌 $\to$ 팝업 $-$0.0765에서 $-$0.0895, 게임 $\to$ 아이돌 $-$0.0752에서
$-$0.0872. 아이돌이 끼지 않은 두 방향은 변화가 없어, 개선이 이 축에서 왔음이
분명하다. 노트 20의 미디어 투입은 추가하니 나빠졌는데 이 축은 좋아졌다 --- 차이는
기능적 정합성이다.
\end{abstract}
]

\section{가장 값진 축이 가장 안 재진다}

노트 22의 결론은 이랬다.

\begin{quote}\fontsize{9.2}{11}\selectfont
매장 노출도가 가장 값지다. 빼면 전이가 43\% 악화된다. 그런데 이 축은 아이돌에서
32\%, 게임에서 0\%다. \emph{가장 값진 축이 가장 안 재진다.}
\end{quote}

그리고 투자 우선순위 1번으로 ``매장 노출도를 다른 도메인에서 관측 가능하게 만든다 ---
이것 하나가 표본을 두 배 늘리는 것보다 낫다''를 적었다. 실행했다.

\section{두 번 실패했다}

\begin{figure}[t]
\centering

\includegraphics[width=\textwidth]{figs/try.pdf}
\caption{세 번의 시도. 앞의 둘은 도메인 독립 물리량을 노렸고 실패했다.}
\end{figure}

노트 20이 미디어 투입에서 배운 것은 ``이름이 같아도 도메인마다 다른 물리량을 재면
인자 공간을 왜곡한다''였다. 그래서 이번에는 물리량부터 정의했다.

\paragraph{매장 노출도 $=$ 경쟁 밀도의 역수}
``그 접점에 우연히 노출될 확률''이 이 축의 물리적 내용이다. 같은 시기에 열린 것이
많을수록 확률이 낮아진다. 세 도메인에 모두 적용 가능한 정의로 보였다.

\begin{itemize}[leftmargin=1.1em,itemsep=0.15em]
\item 팝업 --- 반경 500m 내 동시 운영 팝업 수(이미 보유)
\item 아이돌 --- 데뷔 당월 데뷔 그룹 수
\item 게임 --- 출시 당월 Steam 출시작 수
\end{itemize}

\paragraph{게임에서 실패 --- API가 날짜를 무시한다}
Steam 검색 엔드포인트에 \texttt{released\_date\_range}를 걸어 월별 총수를
받으려 했는데 모든 구간이 같은 값을 돌려줬다(2016년 6월, 2020년 6월, 2024년 6월,
2026년 6월 모두 116,832). 노트 13에서 \texttt{filter=topsellers}와 함께 쓸 때
날짜가 무시된 것과 같은 문제이며, 이번에는 정렬을 바꿔도 해결되지 않았다.
그 값은 ``출시일이 기록된 전체 게임 수''로 보인다.

\paragraph{아이돌에서 실패 --- 표본이 census가 아니다}
아이돌은 내 표본으로 월별 데뷔 수를 셀 수 있다. 세어 봤더니 초동과의 상관이
$+$0.078이고 \emph{부호가 예상과 반대}였다(혼잡한 달의 초동이 오히려 높다).

원인은 표본의 성격이다. 내 274건은 시장 census가 아니라 크롤링으로 모은 것이고,
화제성 높은 시기를 더 많이 잡았다. 그러면 ``월별 데뷔 수''가 경쟁 밀도가 아니라
\emph{그 달이 활발했는가}를 재게 된다.

\claimx{경쟁 밀도로 정의한 매장 노출도는 시장 census를 요구한다. 편향된 표본에서
밀도를 세면 경쟁이 아니라 그 시기의 활발함을 재게 된다.}{같은 기간의 공식 데뷔
집계(가요 관련 협회 통계 등)를 구해 다시 세었을 때 부호가 음으로 돌아서면,
정의는 옳고 표본만 문제였다.}

이것은 왜 이 축이 팝업에서만 100\%인지도 설명한다. 팝업은 우리가 시장 안에
있으므로 상권 데이터를 안다. 밖에서 크롤링으로 보는 도메인에서는 census가 없다.

\section{세 번째에 됐다 --- 라벨 없는 소속사 이력}

경쟁 밀도를 포기하고 다른 경로를 잡았다. 아이돌의 매장 노출도는 원래 ``소속사가
확보하는 유통·음방 자리''로 정의했고, 소속사 사전 데뷔의 \emph{초동 평균}으로
쟀다. 그것이 33\%뿐이었던 이유는 사전작에 초동 라벨이 있어야 했기 때문이다.

\textbf{건수만 세면 라벨이 필요 없다.}

\begin{figure}[t]
\centering
\includegraphics[width=\columnwidth]{figs/mono.pdf}
\caption{소속사 사전 데뷔 건수와 초동. 단조 관계가 뚜렷하다.}
\end{figure}

\begin{table}[t]
\centering\tsize
\caption{소속사 사전 데뷔 건수별 초동(n=71).}
\begin{tabular}{lrr}
\toprule
소속사 이력 & n & 초동 log10 평균 \\
\midrule
첫 데뷔 & 43 & 4.70 \\
1건 경험 & 12 & 5.22 \\
2건 이상 & 16 & 5.34 \\
\bottomrule
\end{tabular}
\end{table}

\begin{table}[t]
\centering\tsize
\caption{두 정의의 비교.}
\begin{tabular}{lrr}
\toprule
정의 & 커버리지 & 초동과 r \\
\midrule
사전작 초동 평균 (기존) & 33\% & $+$0.519 \\
사전 데뷔 건수 (신규) & \textbf{90\%} & $+$0.209 \\
\bottomrule
\end{tabular}
\end{table}

상관은 약해지지만 커버리지가 세 배 가까이 오른다. 둘을 결합했다 --- 사전작 초동이
있으면 그것을 쓰고, 없으면 건수로 잡는다. 최종 커버리지 90\%다.

\section{노트 22의 예측이 맞았다}

\begin{figure}[t]
\centering
\includegraphics[width=\textwidth]{figs/gain.pdf}
\caption{매장 노출도 확보 전후. 별표가 아이돌이 낀 방향이며, 그 셋만 개선됐다.}
\end{figure}

\begin{table}[t]
\centering\tsize
\caption{프로크루스테스 전이, 순열 3,000회.}
\begin{tabular}{lrrrr}
\toprule
& \multicolumn{2}{c}{확보 전 (33\%)} & \multicolumn{2}{c}{확보 후 (90\%)} \\
방향 & 차이 & p & 차이 & p \\
\midrule
팝업 $\to$ 아이돌 & $-$0.0716 & 0.0017 & \textbf{$-$0.0905} & \textbf{0.0003} \\
아이돌 $\to$ 팝업 & $-$0.0765 & 0.0293 & \textbf{$-$0.0895} & \textbf{0.0203} \\
게임 $\to$ 아이돌 & $-$0.0752 & $<$0.0001 & \textbf{$-$0.0872} & \textbf{$<$0.0001} \\
\midrule
게임 $\to$ 팝업 & $-$0.0671 & 0.0010 & $-$0.0671 & 0.0010 \\
팝업 $\to$ 게임 & $-$0.0280 & 0.0183 & $-$0.0280 & 0.0183 \\
아이돌 $\to$ 게임 & $+$0.0019 & 0.4110 & $+$0.0003 & 0.4120 \\
\bottomrule
\end{tabular}
\end{table}

아이돌이 낀 세 방향이 16\%에서 26\% 강해졌다. 아이돌이 끼지 않은 두 방향은
소수점 넷째 자리까지 변화가 없다 --- 개선이 이 축에서 왔음이 분명하다.

\claimx{매장 노출도의 커버리지를 33\%에서 90\%로 올리자 그 도메인이 낀 전이가
모두 16--26\% 강해졌다. 노트 22가 축 절제로 예측한 값어치가 실제 확보에서
확인됐다.}{다른 축을 같은 폭으로 확보했을 때도 비슷한 개선이 나오면, 개선은
이 축의 값어치가 아니라 커버리지 자체의 효과다.}

\section{왜 이 축은 되고 미디어 투입은 안 됐나}

노트 20에서는 게임의 미디어 투입 커버리지를 38\%에서 100\%로 올렸다가 전이가
나빠져 되돌렸다. 여기서는 33\%에서 90\%로 올렸더니 좋아졌다. 같은 종류의 작업인데
결과가 갈린다.

차이는 \textbf{기능적 정합성}이다.

\begin{table}[t]
\centering\tsize
\caption{두 축의 도메인별 측정과 그 기능.}
\begin{tabular}{lll}
\toprule
축 & 도메인 & 재는 것 \\
\midrule
미디어 & 팝업 & 계획된 홍보 투입 \\
투입 & 아이돌 & 사전 노출의 \emph{결과} \\
 & 게임 & 배급 \emph{구조} \\
\midrule
매장 & 팝업 & 상권의 노출 정도 \\
노출도 & 아이돌 & 소속사의 유통 확보력 \\
\bottomrule
\end{tabular}
\end{table}

미디어 투입은 계획과 결과와 구조라는 서로 다른 \emph{종류}를 재고 있었다. 매장
노출도는 팝업과 아이돌 모두 ``접근 경로를 얼마나 확보했는가''를 잰다 --- 측정
도구는 상권 등급과 소속사 이력으로 다르지만 기능이 같다.

\claimx{축을 확보할 때 중요한 것은 측정 도구의 동일성이 아니라 기능의 동일성이다.
같은 기능을 재면 도구가 달라도 되고, 다른 기능을 재면 도구를 통일해도 안 된다.}{기능이
같아 보이는데도 확보 후 전이가 나빠지는 축이 나오면, 이 기준으로는 부족하고
더 엄밀한 정합성 판정이 필요하다.}

\section{남은 것}

게임의 매장 노출도는 여전히 0\%다. Steam은 스토어 내 노출 데이터를 공개하지 않고,
경쟁 밀도는 API 제약으로 못 얻는다. 게임에서 ``접근 경로 확보''에 해당하는 관측
가능한 신호를 아직 찾지 못했다.

그리고 아이돌 $\to$ 게임은 여전히 유의하지 않다(p=0.412). 이 방향은 노트 13부터
한 번도 성공한 적이 없다.

\section{한계}

\begin{itemize}[leftmargin=1.1em,itemsep=0.15em]
\item 사전 데뷔 건수는 \emph{내 표본 안에서} 센다. 소속사가 표본 밖에서 데뷔시킨
      그룹은 세지 못하므로 과소 계상된다.
\item 두 정의를 결합했으므로 33\%는 초동 평균, 나머지 57\%는 건수다. 한 축 안에
      두 눈금이 섞여 있다.
\item 개선폭 16--26\%는 단일 측정이다. 순열 p는 있지만 개선폭 자체의 신뢰구간은 없다.
\item 경쟁 밀도 정의의 실패를 ``표본이 census가 아니어서''로 설명했지만 검정하지
      않았다. 공식 집계를 구해 확인해야 한다.
\end{itemize}

\section{다음}

\begin{enumerate}[leftmargin=1.3em,itemsep=0.15em]
\item 게임의 매장 노출도 --- 접근 경로 확보에 해당하는 관측 가능한 신호 탐색.
\item 한 축 안에 섞인 두 눈금 정리 --- 건수 정의로 통일했을 때의 손실 측정.
\item 아이돌 $\to$ 게임 방향의 실패 원인 --- 스물세 편째 미해결.
\item 개선폭에 신뢰구간 --- 부트스트랩으로 확보 효과 자체를 검정.
\end{enumerate}

\vspace{0.6em}
\hrule height 0.3pt
\vspace{0.5em}
{\footnotesize
\textbf{참고문헌}\\[0.3em]
Meredith, W. (1993). Measurement invariance, factor analysis and factorial
invariance. \emph{Psychometrika}, 58(4), 525--543.\\[0.15em]
Schönemann, P. H. (1966). A generalized solution of the orthogonal Procrustes
problem. \emph{Psychometrika}, 31(1), 1--10.\\[0.15em]
Guyon, I., Elisseeff, A. (2003). An introduction to variable and feature selection.
\emph{JMLR}, 3, 1157--1182.\\[0.15em]
Ben-David, S. et al. (2010). A theory of learning from different domains.
\emph{Machine Learning}, 79(1--2), 151--175.
}

\end{document}
